% concatenated from ON_THE_FUNDAMENTAL_SPECTRAL_GAP_OF_WEIGHTED_SCHR__DINGER_OPERATORS.tex
\documentclass[12pt,a4paper]{amsart}
\usepackage[utf8]{inputenc}
\usepackage{tikz}
\usepackage[english]{babel}
\usepackage[T1]{fontenc}
\usepackage{amsmath}
\usepackage{amsfonts}
\usepackage{amssymb}
\usepackage{graphicx}
\usepackage[left=2.5cm,right=2.5cm,top=3cm,bottom=2.5cm]{geometry}
\usepackage{ulem}
%\usepackage{hyperref}
%\usepackage{amssymba
%%%%%%%%%%%%%%%%%%%%%%%%%%%%%%%%%%%%%%%%%%%%%%%%%%%%%%%%%%%%%%%%%%%%%%%%%%%%%%
%\usepackage{geometry}
%\geometry{top=2cm,bottom=2cm,hscale=0.8,headheight=15pt, textwidth=168mm} % Mise en page des marges etc.
\usepackage{color}
%%%%%%%%%%%%%%%%%%%%%%%%%%%%%%%%%%%%%%%%%%%%%%%%%%%%%%%%%%%%%%%%%%%%%%%%%%%%
\newtheorem{defi}{Definition}[section]
\newtheorem{thm}{Theorem}[section]
\newtheorem{ex}{Example}[section]
%\newtheorem{proof}{Proof}[section]
\newtheorem{rem}{Remark}[section]
\newtheorem{rems}{Remarks}[section]
 \newtheorem{prop}{Proposition}[section]
\newtheorem{lem}{Lemma}[section]
% \maketitle \numberwithin{equation}{section}
\newtheorem{theorem}{Theorem}[section]
\newtheorem{remark}[theorem]{Remark}
\newtheorem{conj}[theorem]{Conjecture}
\newtheorem{opn}[defi]{Open Problem}%[section]
\newtheorem{lemma}[theorem]{Lemma}
\newtheorem{proposition}[theorem]{Proposition}
%\newtheorem{def}{Definition}[section]
\newtheorem{cor}{Corollary}[section]
%\renewcommand{\abstractname}{\bf{Acknowledgements}}
%\newtheorem{cor}{Corollary}[section]
\def\R{{\rm {I  \! R}}}
\def\N{{\rm {I  \! N}}}
\newcommand{\eh}{\textcolor{red}}
\newcommand{\ma}{\textcolor{cyan}}
\usepackage{color}

\begin{document}

 \title [Fundamental spectral gap] {On the fundamental spectral gap of weighted Schr\"{o}dinger operators }


\address{ \newline
Team of Modeling and Scientific Computing, 
Department of Mathematics, 
Multidisciplinary Faculty of Nador, 
University of Mohammed First, Morocco}
\email{m.ahrami@ump.ac.ma}

%\address{ \newline
%Team of Modeling and Scientific Computing, 
%Department of Mathematics, 
%Multidisciplinary Faculty of Nador, 
%University of Mohammed First, Morocco}
%\email{z.elallali@ump.ma}
%%\newline

%\address{School of Mathematics,
%Georgia Institute of Technology,\
%Atlanta GA 30332-0160, USA} 
%\email{harrell@math.gatech.edu}

\subjclass[2010]{34L15, 34B27, 35J60, 35B05.}

\keywords{Fundamental spectral gap, weighted Schr\"{o}dinger operators, convex potential, concave weight, Dirichlet boundary conditions }



%\address{Laboratory of Applied Mathematics and
%Information Systems ,Department of Mathematics and Computer,
%Faculty Multidisciplinary of Nador,University Mohammed Premier, Morocco.}
%\email{elallali@hotmail.com}


\author{Mohammed Ahrami}
 \maketitle

%\address{School of Mathematics,
%Georgia Institute of Technology,\
%Atlanta GA 30332-0160, USA.} 
%\email{ harrell@math.gatech.edu}

 \maketitle

\begin{abstract}
We characterize the potential $V(x)$ that minimizes the fundamental spectral gap of weighted Schr\"{o}dinger operators on the interval $[0,\pi]$ subject to Dirichlet boundary conditions, under the constraint that the potential $V(x)$ is convex and that the weight function $w(x)$ is concave.

\end{abstract}

%\subjclass[2010]{34B27, 35J60, 35B05.}


%{\em Keywords. Fundamental gap spectral, Schr\"odinger operator, singl-well potentials, Dirichlet boundary conditions.}

\section{Introduction}
%\noindent



%\bigskip
%\noindent



In this article, we study the minimization problem of the {\it fundamental eigenvalue gap}
$\Gamma := \lambda_2 - \lambda_1$ of weighted Schr\"{o}dinger operators

\begin{equation}\label{ePnoP}
H(V,w) u :=  -u^{\prime \prime} +V(x) u  = \lambda  w(x) u,
\quad\quad x\in [0,\pi].
\end{equation}
on a finite interval $[0,\pi]$ with Dirichlet boundary conditions, i.e., $u(0) = u(\pi) = 0$.

First, we provide some motivation for studying the fundamental gap and the importance of its bounds. There are both physical and mathematical reasons to study the fundamental gap. For example, in quantum mechanics, the fundamental gap represents
the energy needed to achieve the first excited state from the ground state of the
particle described by the considered Schr\"odinger operators, the so-called excitation
energy. The fundamental gap is also of importance in quantum field theory and 
statistical mechanics. To give an example, Van den Berg in \cite{Berg} applied gap 
results of the Laplace operator to provide sufficient conditions for a free boson gas 
to fill the ground state alone under the thermodynamic limit macroscopically. 

 The optimal estimates of the fundamental spectral gap of Schr\"odinger operators have received considerable attention in the last decades, and have been extensively studied since the 1980s; cf. \cite{Ash1, Ash2,Ash3, Ash4, ahrami, Ahrami3, Horvath1,Horvath2,Lavine} with most prior works assuming that $w(x)=1$.

 In the unweighted case,
where \(w(x)=1\), classical results of Ashbaugh and Benguria \cite{Ash1} established
sharp lower bounds for the fundamental gap under symmetry and single-well
assumptions on the potential. Later, Lavine \cite{Lavine} proved that the constant potential minimizes the fundamental gap among
all convex potentials and satisfies the lower bound


\begin{equation}\label{conjecture1}
 \Gamma\geqslant \frac{3\pi^{2}}{l^{2}} 
\end{equation}

\noindent where $l$ denotes the length of the interval.

This result revealed a remarkable rigidity phenomenon
for Schr\"odinger operators with convex potentials and became one of the central results in one-dimensional
spectral gap theory.

Some time later, in \cite{Andrews}, Andrews and Clutterbuck showed that the
fundamental gap of a Dirichlet Schr{\"o}dinger operator with semi-convex potential $q$ on a bounded convex domain $\Omega\in \mathbb{R}^{d}$ with diameter $D$ satisfies the estimate \eqref{conjecture1}, thus proving the so-called fundamental gap conjecture in essentially full generality. 

For graphs, Ahrami et al. \cite{AHEK} proved an interesting phenomenon: on a metric tree graph the optimal
potential for Schr\"odinger operators with convex potentials must be piecewise affine and there are graphs where the constant potential
does not minimize.

More recently, Ahrami et al. \cite{AHRA} studied the problem (\ref{ePnoP}) for single-well potentials and single-barrier weights and proved that for weighted Schr\"{o}dinger operators, both the optimal potential and the optimal weight are step functions with a specific characterization.

Our goal is to determine
whether the classical phenomenon discovered by Lavine
persists in the weighted setting, namely whether constant potentials continue to minimize the fundamental gap in the presence of nontrivial weight. The presence of a nontrivial weight significantly changes the spectral structure of the problem and introduces new interactions between the geometry
of the potential and the inhomogeneity of the medium.


Our first result shows that any minimizing convex potential can be reduced to an affine potential without increasing the spectral gap. Similarly, every concave weight can be reduced to an affine weight. This reduction is obtained through a perturbative argument based on the Feynman-Hellmann formula together with qualitative properties of the first two eigenfunctions. We then analyze affine weights in greater detail and prove that, under a suitable quantitative condition on the weight,
\[p > 5m\pi,
\]
where \(w(x)=mx+p\), any gap-minimizing convex potential must necessarily be constant.

%The results obtained in this paper provide a weighted analogue of classical spectral gap minimization phenomena and reveal how strong inhomogeneities in the weight dominates the influence of the potential \cite{Filo}. From a physical perspective, this indicates that sufficiently strong spatial inhomogeneity suppresses the effect of nonuniform convex potentials, leading to constant configurations as optimal minimizers of the excitation energy.

Compared with the previous work of Ahrami et al. \cite{AHRA} on single-well potentials and single-barrier functions, the present paper
reveals a fundamentally different optimization mechanism. In the
single-well setting, the minimizers are typically discontinuous step
functions. Furthermore, in this paper the
optimization procedure reduces the problem to affine structures and eventually
recovers constant potentials under suitable quantitative assumptions on the
weight. Thus, the present work establishes a new bridge between the direct
optimization developed in \cite{ElAHar,
AHRA} and Lavine's result \cite{Lavine}.

The paper is organized as follows. In Section 2, we recall several preliminary properties of the fundamental eigenvalue gap, including a Feynman-Hellmann formula for eigenvalue variations and results on the monotonicity of the eigenfunction ratio $\frac{u_2}{u_1}$. In Section 3 we then establish in Corollary \ref{existence} general existence results by invoking compactness of the sets of functions with the properties of interest,
which in turn follow from a general argument using either the Blaschke selection principle. With these results in hand, we turn to characterizing the optimizers in Theorem \ref{step} and finally in Theorem \ref{step2}, we prove that for sufficiently large affine weights, the minimizing convex potential is necessarily constant.  








\section{Preliminaries and basics}
Throughout the paper, we study the weighted Schr\"{o}dinger eigenvalue problem
\begin{equation*}\label{eq}
H(V,w) u :=  -u^{\prime \prime} +V(x) u  = \lambda  w(x) u
\end{equation*}

subject to Dirichlet boundary conditions $u(0)=u(\pi)=0$
with a measurable bounded potential $V:[0,\pi] \longrightarrow
 \mathbb{R}$ and measurable bounded weight $w:[0,\pi] \longrightarrow
 \mathbb{R^+}$.

\noindent Under the above assumptions on $V$ and $w$, the operator $H(V,w)$ is self-adjoint with compact resolvent and therefore possesses a purely discrete spectrum; see, e.g., \cite{AHRA,Wei}. 


%According to \cite{AHRA}, the weighted Schr\"{o}dinger operator $H$ is self-adjoint and has discrete spectrum \cite{Wei}. 

%We start with the basic framework we will use throughout the paper.
\noindent We denote the eigenvalues of $H$ by


%The eigenvalues of $H$ form a strictly
%increasing sequence that accumulates only at infinity.
$$
0<\lambda_1(V, w)<\lambda_2(V, w)<\lambda_3(V, w)<...
$$
indicating the dependence on the potential $V$ and the weight $w$. Furthermore, all eigenvalues are simple; see, e.g., \cite{Zett, Wei}.

% We will abbreviate this notation
%at some points  by simply writing $\lambda_n$ for the eigenvalues.  Likewise the corresponding normalized eigenfunctions $u_n(V,w, x)$ will be written simply as $u_n(x)$.
The main object of interest in this paper is the fundamental gap $\Gamma[V, w]$ defined by  

 $$\Gamma[V, w] = \lambda_2(V, w)- \lambda_1(V, w).$$
It follows immediately from the definition that the fundamental gap of a weighted Schr\"{o}dinger operators subject to
Dirichlet boundary conditions on an interval $[0,\pi]$ is invariant under the addition of a constant, i.e.,
$ \Gamma[V+c,w]= \Gamma[V,w]$.
As a consequence of this remark, which will be very important later on, if we have 
to deal with affine potentials, we only need to consider potentials of the form
$V(x)=ax $  with $a>0$. 
 %It is clear that
%depends on the potential $V$ and the weight $w$ and the length of the underlying interval and  that 
%the fundamental gap $\Gamma(V, w)$ is unaffected by adding a constant, that
%is $$\Gamma(V, w)=\Gamma(V+c, w).$$ 
%to indicate the dependence on the fundamental gap on coefficients in {\eqref{ePnoP}} with respect to which we wish to optimize.  

We recall several standard facts from spectral gap theory; see for instance \cite{Ash1,Lavine,ElAHar}.

%the basic framework that are standard in the literature; see, e.g., \cite{Ash1,Lavine,ElAHar}..

%we review and slightly extend some useful observations about $\Gamma$ that are familiar from previous sources such as \cite{Ash1,Lavine,ElAHar}.
%Most importantly, there is an explicit formula for the first derivative with respect to perturbations of $V$ and $w$:

\begin{lem}\label{F-H}
Suppose that $V(.,\kappa)$ and $w(.,\kappa)$ are one-parameter families of real-valued, locally $L^1$ functions, \( C^1 \) in \( \kappa \) for \( \kappa \) in some real interval, with
$\inf_{x} V(x,\kappa) > - \infty$, $C \ge w(x,\kappa) \ge \frac 1 C$ for some $C > 0$, and
$\frac{\partial V}{\partial \kappa}(x,\kappa)$
and $\frac{\partial w}{\partial \kappa}(x,\kappa) \in L^1(0, \pi)$. Without loss of generality we standardize the first two normalized eigenfunctions $u_{1,2}(x)$ associated with 
$\lambda_{1,2}$
so that
$u_{1,2}(x) > 0$ for $0 < x < \epsilon$ for some $\epsilon$.  Then
\begin{enumerate}
\item \[
\frac{d\lambda_{n}(\kappa)}{d\kappa}=-\lambda_{n}\int_{0}^{\pi}\frac{\partial w}{\partial \kappa}(x,\kappa)u_{n}^{2}(x,\kappa)dx+\int_{0}^{\pi}\frac{\partial V}{\partial \kappa}(x,\kappa)u_{n}^{2}(x,\kappa)dx .
\]
\item
$\dfrac{u_{2}}{u_{1}} $ is decreasing on $(0,\pi)$.

\item The equation $ \vert u_{1}(x)\vert=\vert u_{2}(x)\vert $ has
either one or two solutions on $(0,\pi)$.
 \item There exist two points $x_{-}$ and $x_{+}$,
  $ 0 \leq x_{-}< x_{+} \leq \pi $, at least one of which is interior to $(0, \pi)$, such that 
$u_{1}^{2}(x)>u_{2}^{2}(x)$ on $(x_{-},x_{+})$ and $u_{1}^{2}(x)\le u_{2}^{2}(x)$ on $(x_{-},x_{+})^c$.
\item
The equation $ \lambda_1 \vert u_{1}^2(x)\vert=\lambda_2 \vert u_{2}^2(x)\vert $ has
either one or two
solutions on $(0,\pi)$.
 \item 
There exist two points $\widehat{x}_{-}$ and $\widehat{x}_{+}$,
  $ 0 \leq \widehat{x}_{-}< \widehat{x}_{+} \leq \pi $, at least one of which is interior to $(0, \pi)$, such that 
$\lambda_1 u_{1}^{2}(x)> \lambda_2 u_{2}^{2}(x)$ on $(\widehat{x}_{-},\widehat{x}_{+})$ and $\lambda_1 u_{1}^{2}(x)\le \lambda_2 u_{2}^{2}(x)$ on $(\widehat{x}_{-},\widehat{x}_{+})^c$.
\end{enumerate}

\end{lem}

\noindent
{For the proof,  see \cite{AHRA}.}



\begin{lem}\label{G'''}
Let $G:[0,\pi]\to\mathbb{R}$ be a function such that $G^{''}(x)$ is absolutely continuous. If $u$ satisfies \eqref{ePnoP}, where $V(x),w(x)$ are differentiable, then
\begin{align*}
G(\pi)[u'^{2}(\pi)&+(\lambda w(\pi)-V(\pi))u^{2}(\pi)]-G(0)[u'^{2}(0)+(\lambda w(0)-V(0))u^{2}(0)]\\
&\quad\quad\quad\quad+ \dfrac{1}{2}[G''(\pi)u^{2}(\pi)-G''(0)u^{2}(0)]\\
&=\int_{0}^{\pi}\left[2G'(\lambda w-V)+G(\lambda w'-V')+\frac{G'''}{2}\right]
u^{2}(x)dx.
\end{align*}
\end{lem}


\begin{proof}  
We have
\begin{align*}
G(\pi)[u'^{2}(\pi)&+(\lambda w(\pi)-V(\pi))u^{2}(\pi)]-G(0)[u'^{2}(0)+(\lambda w(0)-V(0))u^{2}(0)]\\
 &=\int_{0}^{\pi}\frac{dG(x)}{dx}[u'^{2}(x)+(\lambda w(x)-V(x))u^{2}(x)]dx  \\
&=\int_{0}^{\pi}G'[u'^{2}+(\lambda w-V)u^{2}]+G[2u'u''+(\lambda w'-V')u^{2}+2(\lambda w-V)uu']dx.\\
&=\int_{0}^{\pi}G'[u'^{2}+(\lambda w-V)u^{2}]+G[-2(\lambda w-V)uu'+(\lambda w'-V')u^{2}+2(\lambda w-V)uu']dx\\
&=\int_{0}^{\pi}G'[u'^{2}+(\lambda w-V)u^{2}]+G(\lambda w'-V')u^{2}dx
\end{align*}
and 
\begin{align*}
0 &=\int_{0}^{\pi}(G'uu')'dx\\
&=\int_{0}^{\pi}(G''uu'+G'u'^{2}+G'uu'')dx\\
&=\int_{0}^{\pi}[G''uu'+G'u'^{2}-G'(\lambda w-V)u^{2}]dx\\
&=\frac{1}{2}[G''(\pi)u^{2}(\pi)-G''(0)u^{2}(0)]-\frac{1}{2}\int_{0}^{\pi}G'''u^{2}dx-\int_{0}^{\pi}G'(\lambda w-V)u^{2}dx+\int_{0}^{\pi}G'u'^{2}dx.
\end{align*}
(Since $G''$ is assumed absolutely continuous $G'''$ is defined a.e. and integration by parts is valid.)  Thus
$$\frac{1}{2}[G''(\pi)u^{2}(\pi)-G''(0)u^{2}(0)]=\int_{0}^{\pi}
\left[\frac{G'''u^{2}}{2}+G'(\lambda w-V)u^{2}-G'u'^{2}\right]dx.
$$
\noindent Combining the previous identities yields the desired formula
 $$ G(\pi)[u'^{2}(\pi)+(\lambda w(\pi)-V(\pi))u^{2}(\pi)]-G(0)[u'^{2}(0)+(\lambda w(0)-V(0))u^{2}(0)]$$
$$ + \dfrac{1}{2}[G''(\pi)u^{2}(\pi)-G''(0)u^{2}(0)]$$
$$=\int_{0}^{\pi}\left[\frac{G'''u^{2}}{2}+G'(\lambda w-V)u^{2}-G'u'^{2}+G(\lambda w'-V')u^{2}+G'((\lambda w-V)u^{2}+u'^{2})\right]dx. $$
\end{proof}


%\begin{proof}  
%Because $C>0$, the weighted Hilbert space is norm-equivalent to the one with $w(x) = 1$, and hence
%$\frac{\partial V}{\partial \kappa}$ and $\frac{\partial w}{\partial \kappa}$ are relatively bounded perturbations, so Kato's theory of 
%analytic perturbations applies.  Since the eigenvalues with separated homogeneous boundary conditions are simple, this 
%justifies the use of a formal expansion to calculate the effect of the perturbation, \`a la Feynman-Hellmann:
%Denoting $\dot{u}=\frac{\partial u}{\partial \kappa}$, differentiation of  {\eqref{ePnoP}} with respect to $\kappa$

%\end{proof} 

%We next adapt the monotonicity argument of \cite{Ash1,ElAHar}
%to incorporate the weight:






\section{Characterization of optimizers}
 In this section, we look at the explicit form of the minimizing potential of the fundamental spectral gap for weighted Schr\"{o}dinger operators under Dirichlet boundary conditions {\eqref{ePnoP}} 
{closely following the strategy in \cite{ElAHar}.}



\subsection{The class of convex potentials and concave weight functions}

 We define here several classes of functions that will play a role below.

\begin{defi}

For $1<M \le \infty$ and $0 < N_{<,>} < \infty $, let
$$ \mathcal{C}_M=\lbrace V(x): 0\le V(x)\le M : V(x) 
\textrm{ is convex on }[0,\pi] \rbrace .$$
and
$$ \mathcal{S}_{N_{<,>}}=\lbrace w(x):  N_< \leq w(x) \leq N_>: w(x) 
\textrm{ is concave on }[0,\pi] \rbrace .$$
\end{defi}



\begin{defi}
Consider the weighted Schr\"{o}dinger operators with Dirichlet boundary conditions
\begin{align}\label{s-l}
    {-u''+V(x)u=\lambda w(x)u}&\\
    u(0)=u(\pi)=0&,\nonumber
\end{align}                                                   
where $V$ is a bounded measurable convex function and $w$ is a bounded measurable concave density.
If there exist $V_{*} \in \mathcal{C}_M$ and 
$w_{*} \in  \mathcal{S}_{N_{<,>}}$ such that
\[
\Gamma(V_{*},w_{*})= \inf \lbrace
  \Gamma(V,w),V \in \mathcal{C}_M,
~w\in \mathcal{S}_{N_{<,>}} \rbrace,
\]
then we call the function $V_{*}$ an {\rm optimal potential} and the function  $w_{*}$ an {\rm optimal density} for problem \eqref{s-l}.
\end{defi}


\begin{conj}
We consider the weighted Schr\"{o}dinger operators

\begin{equation}\label{eq12}
\left\{
\begin{array}{lr}
H(V,w) u :=  -u^{\prime \prime} +V(x) u  = \lambda  w(x) u
 \\
u(0)=u(\pi)=0 & 
\end{array}
\right.
\end{equation}
For a given bounded measurable, strictly positive weight function $w$, any 
gap-minimizing
optimal potential $V_{*}\in \mathcal{C}_M $ is constant.
\end{conj}


\begin{rem}\label{rem1}
The assumptions of boundedness and measurability ensure that the weighted Schr\"{o}dinger operators $H(V,w)$  is self-adjoint with a discrete spectrum, as discussed in \cite{Wei}.
\end{rem}

{As in \cite{ElAHar}, we will use compactness of the sets $\mathcal{C}_M$ and $\mathcal{S}_{N_{<,>}}$ , following from a Blaschke Selection principle: }

\begin{prop}\label{compact}
For any sequence $f_{n}\in\Lambda$, where $\Lambda=\mathcal{C}_M$ or
respectively $\mathcal{S}_{N_{<,>}}$ with any fixed positive $M, N_{<,>}$, there exist a subsequence $f_{n_{k}}$ and a function $f_{*}$ such that  $f_{n_{k}}(x)\longrightarrow f_{*}(x) \in\Lambda $ for a.e. $x$ with $f_{*} \in \mathcal{C}_M$ or respectively $\mathcal{S}_{N_{<,>}}$.
\end{prop}

\noindent
{For the proof,  see \cite{ElAHar},
Proposition $2.1$.}

\begin{cor}\label{existence}
There exist a potential $V_{*}\in \mathcal{C}_M $ and a density $w_{*}\in \mathcal{S}_{N_{<,>}}$ that minimize $\Gamma[V,w]$.
\end{cor}

\begin{proof}  
The argument is identical to the proof of Corollary 3.1 in \cite{AHRA}, except that, as mentioned, one replaces the Helly compactness theorem with the Blaschke selection principle.
\end{proof}  

We can now seek the explicit form of optimal potential  $V_{*}\in \mathcal{C}_M $ of the problem  
{\eqref{s-l}}. We follow the strategy of \cite{AHRA}.



\begin{thm}\label{step}
For a given bounded measurable, strictly positive weight function $w$, and for every convex 
and non-affine potential $V$ (i.e. not of the form $ax+b$), there exists a linear potential $V_{a}=ax$ such that
$$
 \Gamma[V,w] \geq  \Gamma[V_{a},w].$$
Alternatively, for a given  
bounded measurable potential function $V$, and for every concave and not-affine
weight $w$ (i.e. not of the form $mx+p$), there exists an affine weight $w_m(x)=mx+p$ such that 
$$
 \Gamma[V,w] \geq  \Gamma[V,w_m].$$  
\end{thm}

\begin{proof}  

We assume that $w$ is fixed. 
%Our objective is to demonstrate that all optimal $V_{*}$ must be of the form $ax.$
Let $V$ be a convex non-affine potential and $L_{V}(x)=ax+b$ the affine potential 
such that $V(x_{\pm})=L_{V}(x_{\pm})$.


\noindent By Lemma \ref{F-H} there exist $x_{\pm}$: $0\leqslant x_{-}<x_{+}\leqslant \pi$, for which 
\begin{align*}
    &u_{2}^{2}(x)>u_{1}^{2}(x) \textrm{ on } (0,x_{-})\cup (x_{+},\pi)\quad\textrm{(one of these intervals may be vacuous)}\\
    &u_{1}^{2}(x)>u_{2}^{2}(x) \textrm{ on } (x_{-},x_{+}).
\end{align*}
 
\noindent By the convexity of $V$,
\begin{gather*}
 V-L_{V}\geq 0 \quad \text{on } (0,x_{-})\cup (x_{+},\pi), \\
 V-L_{V}\leq 0 \quad \text{on } (x_{-},x_{+}). 
\end{gather*}
Hence 
$$
\int_0^{\pi}(V(x)-L_{V}(x))(u_2^2(x)-u_1^2(x))dx >0.
$$
Let
\[
 L_{V}(\theta)=\theta V+(1-\theta)L_{V}\quad\text{for }\theta\in(0,1).
\]
 Therefore $\dot{L_{V}}(\theta)=V-L_{V}$.
Then by the Feynman-Hellman formula,
$$
\frac{d\Gamma(L_{V}(\theta))}{d\theta}=
\int_0^{\pi}(V(x)-L_{V}(x))(u_2^2(x)-u_1^2(x))dx >0.
$$
Integrating this inequality with respect to $\theta$ over $(0,1)$, we find that
$$
\int_0^{1}\frac{d\Gamma(L_{V}(\theta))}{d\theta}=\Gamma[V,w]-\Gamma[L_{V},w]\geq 0.
$$
Then $$ \Gamma[V,w]\geq \Gamma[L_{V},w]=\Gamma[ax,w].$$

We now prove the second part of the statement by a similar argument applied to $w$ for fixed $V$.
Let $w$ be a concave, not-affine weight, and $w_{m}=mx+p$ be the affine weight such that $w(\widehat{x}_{\pm})=w_{m}(\widehat{x}_{\pm})$.


Applying Lemma \ref{F-H} once again, we find that 
there exist $\widehat{x}_{\pm}$: $0\leqslant \widehat{x}_{-}<\widehat{x}_{+}\leqslant \pi$, satisfying 
\begin{align*}
    &\lambda_{2} u_{2}^{2}(x)>\lambda_{1} u_{1}^{2}(x) \textrm{ on } (0,\widehat{x}_{-})\cup (\widehat{x}_{+},\pi)\\
    &\lambda_{1}u_{1}^{2}(x)>\lambda_{2}u_{2}^{2}(x) \textrm{ on } (\widehat{x}_{-},\widehat{x}_{+}).
\end{align*}
\noindent By the concavity of $w$  
  $$\displaystyle  \left\{%\left(
    \begin{array}{ll}
   
   w-w_{m} \leqslant  0 \qquad\text{on}~~(0,\widehat{x}_{-})\cup (\widehat{x}_{+},\pi),
   \\
    w-w_{m}\geqslant 0 \qquad\text{on}~~(\widehat{x}_{-},\widehat{x}_{+}).                                                                                      \end{array}                                                                                \right.$$
Therefore
$$
\int_{0}^{\pi}(w-w_{m})(\lambda_2 u_{2}^{2}-\lambda_1 u_{1}^{2})dx <0.
$$ Let 
 $$ L_w(\theta)=\theta w+(1-\theta)w_{m} , ~~~ \theta\in(0,1).$$\\
Consequently
 $$\dot{L_{w}}(\theta)=w-w_{m}.  $$
 Then by the Feynman-Hellman formula,
$$\frac{d\Gamma(L_w(\theta))}{d\theta}=(\lambda_1-\lambda_2)
\int_{0}^{\pi}(w-w_{m})(\lambda_2 u_{2}^{2}-\lambda_1 u_{1}^{2})dx >0.
$$

\noindent Integrating this inequality with respect to $\theta$ over $(0,1)$, we find

$$
\int_{0}^{1}\dfrac{d\Gamma(L_w(\theta))}{d\theta}=\Gamma[V,w]-\Gamma[V,w_{m}]\geqslant 0.
$$
Consequently
$$
\Gamma[V,w]\geqslant \Gamma[V,w_m].
$$

\end{proof}

 In the final part of this section, the aim is to analyze the special case of the optimal potential of the form $V(x)=ax$ with $a\geq0.$ 

\noindent  In the context of single-well potentials and single-barrier weights, it was shown in \cite{AHRA} that optimizers have the following
characterization:
\begin{enumerate}
\item
$V_{*}(x) = 0$ a.e. on a
connected component of $\{x: u_2^2(x) > u_1^2(x)\}$, where
$u_{1,2}$ are the eigenfunctions associated with 
$\lambda_{1,2}$, while on the complement of that interval, 
$V_{*}(x) = M$ 
a.e.  
\item
$w_{*}(x) = N_>$ a.e. on a
connected component of $\{x: \lambda_2 u_2^2(x) > \lambda_1 u_1^2(x)\}$ and on the complement of that interval there exists a constant $\mathcal{M}$ such that $w_{*}(x) = \mathcal{M}$ where $N_< \leq \mathcal{M}<N_>$. 
\end{enumerate}
If $w(x)=1$, Lavine’s \cite{Lavine} result can be reformulated equivalently for weighted Schr\"{o}dinger operators as follows: $$
\Gamma[V,1]\geqslant \Gamma[0,1].
$$
%In fact, we expect the property of the constant potential minimizing the fundamental
%gap $\Gamma[V,w]$, and we will present
 %below this conjecture.


\begin{lemma}\label{lemm}
    Let the first two normalized eigenfunctions $u_{1,2}(x)$ associated with the first eigenvalues $\lambda_{1,2}$ of problem (\ref{s-l}). Then
    \[
\int_0^{\pi}x^2(u_2^2(x) -u^2_1(x))dx\leq\frac{2\pi^2}{\min_{x} w(x)}.
\]


\end{lemma}

\begin{proof}
\noindent Let $w \in  \mathcal{S}_{N_{<,>}}$, since $w$ is concave on the compact interval $[0,\pi]$, it follows that $w$ is continuous and hence attains its minimum. In addition, the strict positivity of $w$ yields \[ \min_{x\in [0,\pi]} w(x)>0. \]
\noindent  Let the first two normalized eigenfunctions $u_{1,2}(x)$ associated with the first eigenvalues $\lambda_{1,2}$ of problem (\ref{s-l}), 
we have $\int_0^\pi w(x)u_n^2(x)dx=1$. Hence 
\[
\int_0^{\pi}u_n^2(x)dx\leq
\frac{1}{\min_{x} w(x)}\int_0^{\pi}u_n^2(x)w(x
)dx=\frac{1}{\min_{x} w(x)}.
\]
Therefore 

\[
\int_0^{\pi}\mid u_2^2(x)-u_1^2(x)\mid dx\leq
\int_0^{\pi}( u_2^2(x)+u_1^2(x)) dx\leq
\frac{2}{\min_{x} w(x)}.
\]
Since 
\[
\left| \int_0^{\pi}x^2 (u_2^2(x)-u_1^2(x))dx \right| \leq
\frac{2\pi^2}{\min_{x} w(x)}.
\]
Then
    \[
\int_0^{\pi}x^2(u_2^2(x) -u^2_1(x))dx\leq\frac{2\pi^2}{\min_{x} w(x)}.
\]
    
\end{proof}
\begin{rem}\label{rem12}
    In particular, for the affine weight $w(x)=mx+p$, $m>0$, in the interval $[0,\pi]$, the function $w$ increases since $w^{\prime}(x)=m>0$.
Therefore $\min_{x} w(x)=w(0)=p.$ Then

\[
\int_0^{\pi}x^2(u_2^2(x) -u^2_1(x))dx\leq\frac{2\pi^2}{p}.
\]
\end{rem}

\begin{thm}\label{step2}
Let $w(x)=mx+p$ be a fixed positive affine weight function with $m>0$. Assume that
\[
p>5m\pi
%\int_0^{\pi}x^2(u_2^2(x) -u^2_1(x))dx<\frac{2\pi}{5m}.
\]

\noindent Then any minimizer of the fundamental gap
\[
\Gamma[V,w] = \lambda_2(V,w) - \lambda_1(V,w)
\]
over the class of convex potentials $ \mathcal{C}_M$ must be constant.


%For affine potentials $V(x)=ax$ and for affine positive weights $w(x)=mx+p$ with $m\ne0$, and we assume that $\Gamma[V,w]<\frac{a}{m}$, then any 
%gap-minimizing
%optimal potential $V_{*} \in \mathcal{C}_M  $ is constant.
 % Alternatively, for a given  
%bounded measurable potential function $V$, the 
%gap-minimizing
%optimal weight $w_*\in \mathcal{S}_{N_{<,>}}$ is constant $\text{if } \lambda< \frac{t}{m}$.
    
\end{thm}
\begin{proof}  
By Theorem \ref{step}, for any convex potential $V$, there exists an affine potential
\[
V_a(x) = ax 
\]
such that
\[
\Gamma[V,w] \geq \Gamma[V_a,w].
\]

%Thus, it suffices to consider affine potentials.

\noindent Let the first two normalized eigenfunctions $u_{1,2}(x)$ associated with the first eigenvalues $\lambda_{1,2}$ of problem (\ref{s-l}). 

\noindent We consider a family of potentials $V^{\theta}=\theta ax$. 
 From Lemma \ref{F-H}, we have
 

\begin{align*}
&\frac{d(\lambda_2(V^{\theta})-\lambda_1(V^{\theta}))}{d\theta}\\
&=\int_0^{\pi} \frac{\partial V^{\theta}}{\partial \theta}[u_2^2(x)-u_1^2(x)]dx\\
 &=a\int_0^{\pi}x[u_2^2(x)-u_1^2(x)]dx.
\end{align*}
 

\noindent Suppose that a critical point of the gap occurs
at some $a>0$, then
$$
\int_0^{\pi}x(u_2^2(x)-u_1^2(x))dx =0.
$$
\noindent Applying Lemma \ref{G'''} for $G(x)=x$, we have
\begin{align*}
&\pi[(u_n^{\prime}(\pi))^{2}+(\lambda_n w(\pi)-a\pi )u_n^{2}(\pi)]\\
 &=\int_0^\pi [2(\lambda_n w-ax)+x(\lambda_n w^{\prime}-a)]u_n^2
(x)dx.
\end{align*}

\noindent Since $u_n$ satisfies the Dirichlet boundary conditions $u_n(0)=u_n(\pi)=0$, 

$$\pi(u_n^{\prime}(\pi))^{2}
 =\int_0^\pi [2(\lambda_n w-ax)+x(\lambda_n w^{\prime}-a)]u_n^2
(x)dx.
$$ 

\noindent Similarly for $G(x)=x^2$ and from the Dirichlet boundary conditions, we find that
$$\pi^2(u_n^{\prime}(\pi))^{2}\\
=\int_0^\pi [4x(\lambda_n w-ax)+x^2(\lambda_n w^{\prime}-a)]u_n^2(x)dx.$$
\noindent Consequently
\begin{align*}
&\int_0^\pi [2(\lambda_1 w-ax)+x(\lambda_1 w^{\prime}-a)]u_1^2
(x)dx\\
&=\frac{1}{\pi}\int_0^\pi [4x(\lambda_1 w-ax)+x^2(\lambda_1 w^{\prime}-a)]u_1^2(x)dx
\end{align*}
\noindent and
\begin{align*}
&\int_0^\pi [2(\lambda_2 w-ax)+x(\lambda_2 w^{\prime}-a)]u_2^2
(x)dx\\
&=\frac{1}{\pi}\int_0^\pi [4x(\lambda_2 w-ax)+x^2(\lambda_2 w^{\prime}-a)]u_2^2(x)dx
\end{align*}





%\noindent By theorem \ref{step}, we have $$
 %\Gamma[V,w] \geq  \Gamma[ax,mx+p].$$
\noindent  Subtracting the two identities corresponding to $n=1$ and $n=2$, and using $
\int_0^{\pi}x(u_2^2(x)-u_1^2(x))dx =0$, we obtain

%\noindent We note that $\int_0^\pi w(x)u_n^2(x)dx=1$, \noindent since $
%\int_0^{\pi}x(u_2^2(x)-u_1^2(x))dx =0$.
%\noindent Combining the previous identities
\begin{align*}
&  
\frac{2\pi}{5}(\lambda_2-\lambda_1)\\
&=\left(m(\lambda_2-\lambda_1) -a\right)\int_0^{\pi}x^2(u_2^2(x) -u^2_1(x))dx.\\
\end{align*}

\noindent Hence
\begin{align*}
&  
(\lambda_2-\lambda_1)  \left(\frac{2\pi}{5}-m\int_0^{\pi}x^2(u_2^2(x) -u^2_1(x))dx\right)\\
&=-a\int_0^{\pi}x^2(u_2^2(x) -u^2_1(x))dx.\\
\end{align*}


\noindent Let
\[
\phi(x) = x^2 - Ax - B
\]
\noindent chosen such that
\[
\phi(x_-) = \phi(x_+) = 0.
\]

\noindent Then


$$\displaystyle  \left\{%\left( 
    \begin{array}{ll}
   
   \phi(x) \leqslant  0 \qquad\text{on}~~(x_-,x_+)^c,
   \\
   \phi(x)\geqslant 0 \qquad\text{on}~~(x_-,x_+).                                                                                      \end{array}                                                                                \right.$$

\noindent We obtain
\[
\int_0^\pi x^2 (u_2^2 - u_1^2)\,dx
= \int_0^\pi \phi(x)(u_2^2 - u_1^2)\,dx
+ A \int_0^\pi x(u_2^2 - u_1^2)\,dx
+ B \int_0^\pi (u_2^2 - u_1^2)\,dx>0.
\]

\noindent By Remark \ref{rem12} \[
\int_0^{\pi}x^2(u_2^2(x) -u^2_1(x))dx\leq\frac{2\pi^2}{p}.
\]
It follows that \[\frac{2\pi}{5m}
\leq 
 \frac{2\pi^2}{p}.
\]
But the assumption
\[
p>5m\pi
\]
which yields a contradiction, since the left-hand side is positive. Then the gap-minimizing
optimal potential $V_{*} \in \mathcal{C}_M $ is constant.
\end{proof}  

%\begin{rem}
    %Theorem \ref{step2} admits a natural interpretation within the framework of weighted quantum systems, where the weight function models a spatially inhomogeneous medium. The result shows that, under 
    % the quantitative condition

%  It would be interesting to determine whether the quantitative condition  
%\[
%p>5m\pi,
%\]
%is optimal or merely technical.

%the influence of nonuniform convex perturbations of the potential is no longer sufficient to decrease the excitation energy, and the optimal configuration is achieved by a constant potential.
%\cite{Filo}.
%\end{rem}


\subsection{Discussion and Open Problems}

In the present paper, we extended several classical results concerning the minimization of the fundamental spectral gap for Schr\"odinger operators to the weighted setting. More precisely, we proved that every minimizing convex potential can be reduced to an affine potential, while every minimizing concave weight can be reduced to an affine weight. Furthermore, under the quantitative condition
\[
p>5m\pi,
\]
for affine weights of the form
\[
w(x)=mx+p,
\]
we established that any minimizing convex potential must necessarily be constant.

These results indicate that the rigidity phenomenon discovered by Lavine in the unweighted case persists, at least partially, in the weighted framework. Nevertheless, several important questions remain open.

First, it is not clear whether the condition
\[
p>5m\pi
\]
is optimal or merely technical. The proof relies on an integral estimate involving the first two eigenfunctions, and it is conceivable that sharper estimates may improve this condition or possibly remove it altogether. This naturally leads to the following question:

\begin{opn}    
Let $w$ be a positive concave weight on $[0,\pi]$. Does the constant potential minimzer the fundamental spectral gap over the class of convex potentials. Equivalently, 
$$ \Gamma[V,w]\geq \Gamma[0,w]
$$
for every convex potential $V$.
%A positive answer would provide a complete weighted analogue of Lavine's theorem
\end{opn}





\subsection*{Acknowledgements} The author would like to thank Joachim Kerner (FernUniversität in Hagen) and James B. Kennedy (Universidade de Aveiro) for reading the manuscript and useful remarks. 




\subsection*{Funding} This research was not supported by any sponsor or funder.
%From the eigenvalue equation:
%\[
%-u'' + axu = \lambda (mx+p)u,
%\]

%one derives (after %between $n=2$ and $n=1$):
%\[
%(m\Gamma - a)\int_0^\pi x^2 (u_2^2 - u_1^2)\,dx = C,
%\]
%for some constant $C > 0$.

%\subsubsection*{Step 10: Contradiction}

%From (3), the integral is positive.  



%We finish the proof by noting that the Cartesian product of the sets $SW_{[0,\pi],M}$
 %and $SB_{[0,\pi],N<,N_>}$ is compact, so a jointly minimizing pair $(V_{*},w_{*})$ exists. The same argument by contradiction then characterizes $(V_{*},w_{*})$ as a pair of step functions in the same way as when one of them is fixed in advance.




%\noindent
%Note: By Corollary \ref{existence}, the existence of minimizers for the individual optimization problems was already guaranteed. The argument in the proof refines this result by characterizing the structure of these minimizers, concluding that both the optimal potential and the optimal weight must be step functions. Finally, since the Cartesian product is compact, the same compactness argument implies the existence of a jointly minimizing pair $(V_{*},w_{*})$, whose components retain the step-function form established earlier.         



\begin{thebibliography}{00}

\bibitem{AHRA}
\newblock{Ahrami, M., El Allali, Z., and Harrell, E. M.,
\em On the fundamental eigenvalue gap of Sturm–Liouville operators,}
\newblock Arch. Math. 126 (2026), 187–197.


\bibitem{ahrami}
\newblock { M. Ahrami and Z. El Allali,
\em Lower bounds on the fundamental spectral gap with Robin boundary conditions,}
\newblock  2021 UNC Greensboro PDE Conference. Electron. J. Diff. Eqns. Conf. 26 (2022), pp 1-11. 

\bibitem{Ahrami}
\newblock {Ahrami, M., Zakaria El Allali and Jamal Ounejma (2025) 
\em “Lower bounds on the eigenvalue ratio with Robin boundary conditions”,}
\newblock Gulf Journal of Mathematics, 19(2), pp. 489-499. 
 
\bibitem{AHEK} 
\newblock {Ahrami, M., El Allali, Z., Harrell, E. M., and Kennedy, J. B., 
\em Optimizing the fundamental eigenvalue gap of quantum graphs.} 
\newblock J. Phys. A: Math. Theor. 57 (2024), 385205.


 
 
\bibitem{Ahrami2}  
\newblock {Ahrami, M. and El Allali, Z., 
\em Optimizing the first eigenvalue of nonlinear quantum graphs.}
\newblock Gulf Journal of Mathematics, 17 (2024), 29-43.

\bibitem{Ahrami3}  
\newblock {Ahrami, M., Zakaria El Allali and Jamal Ounejma 
\em Minimizing the Eigenvalue Ratio for the p-Laplacian Operator}.
\newblock Journal of Nonlinear Modeling and Analysis, 7(6), 2213-2223. 


\bibitem{Andrews}
\newblock {B. Andrews and J. Clutterbuck,}
\newblock {\em Proof of the fundamental gap conjecture,}
\newblock J.\ Amer.\ Math.\ Soc.\ 24 (2011), 899--916.



\bibitem{Ash1}
\newblock {M. Ashbaugh and R. Benguria, 
\em Optimal lower bound for the gap between the first two eigenvalues of one-dimensional  Schr\"odinger operators with symmetric single-well potentials, }
\newblock Proc. Amer. Math. Soc. 105 (1989), 419--424.

\bibitem{Ash2}
\newblock {M. S. Ashbaugh and R. Svirsky, 
\em  Periodic potentials with minimal energy bands, }
\newblock Proc. Amer. Math. Soc., 114 (1992) 69--77.

\bibitem{Ash3}
\newblock {M. Ashbaugh and R. Benguria, 
\em Optimal bounds for ratios of eigenvalues of one-dimensional Schr\"odinger operators with Dirichlet boundary conditions and positive potentials,}
\newblock Comm. Math. Phys., 124 (1989), 403--415.

\bibitem{Ash4}
\newblock {M.~S. Ashbaugh, E.~M. Harrell II, and R. Svirsky,
\em On minimal and maximal eigenvalues gaps and their causes, }
\newblock Pac. J. Math. 147 (1991) 1--24.

\bibitem{Cheng1}
\newblock {Y. H. Cheng, S. Y. Kung, C. K. Law and W. C. Lian, 
\em The dual eigenvalue problems for the Sturm-Liouville system, }
\newblock Computers and Mathematics with Applications, 60 (2010) 2556--2563.



\bibitem{AsHa82}
\newblock {M.~S. Ashbaugh and E.~M. Harrell II,
\em Perturbation theory for shape resonances and large
barrier potentials,}
\newblock Commun. Math. Phys. 83 (1982) 151--170.

%\bibitem{Bognar}
%\newblock {G. Bogn\`{a}r and O. Dosly,
% \em The ratio of eigenvalues of the Dirichlet eigenvalue problem for equations with one-dimensional p-Laplacian,}
%\newblock Abstract and Applied Analyis, 2010 (2010).
\bibitem{Berg}
\newblock{ M. van den Berg,
 \em On condensation in the free-Boson gas and the spectrum of the
Laplacian},
 Journal of  Statistical  Physics 31 (1983), 623--637. 

 
  \bibitem{Doob}
\newblock {J. L. Doob,
\em Measure theory,}
\newblock Graduate Texts in Mathematics (143), 1994 Springer.

\bibitem{ElAHar}
\newblock {Z. El Allali  and E.~M. Harrell II ,
{\em Optimal bounds on the fundamental spectral gap with single-well potentials}, }
\newblock Proc. Amer. Math. Soc., 150,(2022), 57--87.





\bibitem{Bognar}
\newblock {G. Bogn\`{a}r and O. Dosly,
 \em The ratio of eigenvalues of the Dirichlet eigenvalue problem for equations with one-dimensional p-Laplacian,}
\newblock Abstract and Applied Analysis, 2010 (2010).
 
 \bibitem{Kato}
\newblock {T. Kato,
\em Perturbation theory for linear operator,}
\newblock  1980 Springer-Verlag.

\bibitem{Filo}
\newblock {M. Filoche and S. Mayboroda,
\em Universal mechanism for Anderson and weak localization,}
\newblock Proc. Natl. Acad. Sci. USA 109, (2012), no 37, 14761-12766.

\bibitem{Huang1}
\newblock {M. J. Huang,
\em On the eigenvalue ratio for vibrating strings,}
\newblock Proc. Amer. Math. Soc., 127 (1999) 1805--1813.



\bibitem{Horvath1}
\newblock {M. Horv\'ath,
\em On the first two eigenvalues of Sturm-Liouville operators,}
\newblock Proc. Amer. Math. Soc. 131 4 (2002) 1215--1224.
 
\bibitem{Horvath2}
\newblock {M. Horv\'ath and M. Kiss,
\em A bound for ratios of eigenvalues of Schr\"{o}dinger operators with single-well potentials,}
\newblock Proc. Amer. Math. Soc., 134 (2006) 1425--1434.

\bibitem{Huang2}
\newblock {M. J. Huang,
\em The eigenvalue gap for vibrating strings with symmetric densities, }
\newblock Acta Math. Hungar., 117 (2007) 341--348.

\bibitem{Huang3}
\newblock {M. J. Huang,
\em A note on the eigenvalues ratio of vibrating strings, }
\newblock Acta Math. Hungar., 123 (2009) 265--271.

\bibitem{Huang4}
\newblock {M. J. Huang and C. K. Law,
\em Eigenvalue ratios for the regular Sturm-Liouville system,  }
\newblock Proc. Amer. Math. Soc., 124 (1996) 1427--1436.

 \bibitem{Lavine}
 \newblock {R. Lavine,
 \em The eigenvalue gap for one-dimensional convex potentials,}
 \newblock Proceedings of the American Mathematical Society 1994; 121, 
815--821. 
 
 
 \bibitem{Simon}
\newblock {Reed, M. and Simon, B.
\em (1975) Methods of Modern Mathematical Physics: Fourier Analysis, Self-Adjointness. }
\newblock  Vol. 2, Academic Press, New York.
 
 


 
\bibitem{svir2}
\newblock {R. Svirsky,
\em Maximally resonant potentials subject to p-Norm constraints,}
\newblock Pac. J. Math. 129, 357--374  (1987).

\bibitem{YuYa}
\newblock {X. J. Yu and C. F. Yang,
\em The gap between the first two eigenvalues of Schr\"odinger operators with single-well potential,}
\newblock Appl. Math. Comp. 268 (2015) 275--283.

 \bibitem{Wei}
J. Weidmann,
\newblock {\em Linear operators in {H}ilbert spaces}, volume~68 of {\em
  Graduate Texts in Mathematics}.
\newblock Springer-Verlag, New York-Berlin, 1980.
\newblock Translated from the German by Joseph Sz{\"u}cs.

 \bibitem{Garrett}
Garrett Birkhoff and Gian-Carlo Rota,
\newblock {\em  
Ordinary Differential Equations, 4th Edition.  
}
\newblock New York: Wiley, 1989.  See Chapter 10, section 9.
%\newblock Translated from the German by Joseph Sz{\"u}cs.



\bibitem{Zett}
A. Zettl
\newblock {\em Sturm-Liouville theory},
\newblock American Mathematical Society, 2005.






% \bibitem{Yu1}
%\newblock {X. J. Yu and C. F. Yang,
%\em The gap between the first two eigenvalues of Schr\"odinger operators with single-well potential,}
%\newblock Applied Mathematics and Computation 268 (2015) 275--283.
 
\end{thebibliography}
\end{document}




\begin{enumerate}
\item
Minimize $\Gamma(p, V_0 + V_1)$ for $V_1 \in \mathcal{S}_M$ for fixed $M < \infty$.
\item
Minimize $\Gamma(p, V_0 + V_1)$ for $V_1 \in \mathcal{S}_{M,a}$ for fixed $M < \infty$.
\item
Minimize $\Gamma(p, V_0 + V_1)$ for $V_1 \in \mathcal{S}_{M,a,b}$ for fixed $M < \infty$.
\item
Minimize $\Gamma(p, V_0 + V_1)$ for $V_1 \in \mathcal{C}_M$ for fixed $M < \infty$.
\item
Minimize $\Gamma(p, V_0 + V_1)$ for $V_1 \in \mathcal{C}_{M,a}$ for fixed $M < \infty$.
\item
Minimize $\Gamma(p, V_0 + V_1)$ for $V_1 \in \mathcal{C}_{M,a,b}$ for fixed $M < \infty$.
\item
Minimize $\Gamma(p, V_0 - V_1)$ for $V_1 \in \mathcal{S}_M$ for fixed $M < \infty$.
\item
Minimize $\Gamma(p, V_0 - V_1)$ for $V_1 \in \mathcal{S}_{M,a}$ for fixed $M < \infty$.
\item
Minimize $\Gamma(p, V_0 - V_1)$ for $V_1 \in \mathcal{S}_{M,a,b}$ for fixed $M < \infty$.
\item
Minimize $\Gamma(p, V_0 - V_1)$ for $V_1 \in \mathcal{C}_M$ for fixed $M < \infty$.
\item
Minimize $\Gamma(p, V_0 - V_1)$ for $V_1 \in \mathcal{C}_{M,a}$ for fixed $M < \infty$.
\item
Minimize $\Gamma(p, V_0 - V_1)$ for $V_1 \in \mathcal{C}_{M,a,b}$ for fixed $M < \infty$.
\item
Maximize $\Gamma(p, V_0 + V_1)$ for $V_1 \in \mathcal{S}_M$ for fixed $M < \infty$.
\item
Maximize $\Gamma(p, V_0 + V_1)$ for $V_1 \in \mathcal{S}_{M,a}$ for fixed $M < \infty$.
\item
Maximize $\Gamma(p, V_0 + V_1)$ for $V_1 \in \mathcal{S}_{M,a,b}$ for fixed $M < \infty$.
\item
Maximize $\Gamma(p, V_0 + V_1)$ for $V_1 \in \mathcal{C}_M$ for fixed $M < \infty$.
\item
Maximize $\Gamma(p, V_0 + V_1)$ for $V_1 \in \mathcal{C}_{M,a}$ for fixed $M < \infty$.
\item
Maximize $\Gamma(p, V_0 + V_1)$ for $V_1 \in \mathcal{C}_{M,a,b}$ for fixed $M < \infty$.
\item
Maximize $\Gamma(p, V_0 - V_1)$ for $V_1 \in \mathcal{S}_M$ for fixed $M < \infty$.
\item
Maximize $\Gamma(p, V_0 - V_1)$ for $V_1 \in \mathcal{S}_{M,a}$ for fixed $M < \infty$.
\item
Maximize $\Gamma(p, V_0 - V_1)$ for $V_1 \in \mathcal{S}_{M,a,b}$ for fixed $M < \infty$.
\item
Maximize $\Gamma(p, V_0 - V_1)$ for $V_1 \in \mathcal{C}_M$ for fixed $M < \infty$.
\item
Maximize $\Gamma(p, V_0 - V_1)$ for $V_1 \in \mathcal{C}_{M,a}$ for fixed $M < \infty$.
\item
Maximize $\Gamma(p, V_0 - V_1)$ for $V_1 \in \mathcal{C}_{M,a,b}$ for fixed $M < \infty$.
\end{enumerate}
